\documentclass[dvipsnames]{article}
\usepackage[english]{babel}
\usepackage{multirow}
\usepackage[letterpaper,top=2cm,bottom=2cm,left=3cm,right=3cm,marginparwidth=1.75cm]{geometry}
\usepackage{amsmath}
\usepackage{graphicx}
\usepackage{minted}
\usepackage{amsthm}
\usepackage{mathtools}
\usepackage[colorlinks=true, allcolors=blue]{hyperref}
\usepackage{fancyhdr}
\usepackage{enumerate}

\begin{document}
\section*{Week 4. Problem set (solutions by Karim Khabibrakhmanov)}

\thispagestyle{fancy}
\lhead{Karim Khabibrakhmanov (\texttt{k.khabibrakhmanov@innopolis.university}), group B24-DSAI-05}

\begin{enumerate}

  \item Compute asymptotic worst case time complexity of \texttt{solve} in terms of the size $n$ of the input list:
  \begin{enumerate}[(i)]
    \item Express the running time of the \texttt{solve} procedure $T(n)$ as a recurrence relation.
    \item Find the asymptotic complexity of $T(n)$ using the master theorem.
    \item Specify which case of the master theorem applies (if any).
    \item Write down conditions for this case that need to be checked (write down specialized version for a particular recurrence).
    \item Prove that the conditions are satisfied.
    For asymptotic notation, each property used in a proof \textbf{must} be either
    proven explicitly or properly referenced (e.g. by citing \cite{Cormen2022} or a particular Lecture/Tutorial slide).
  \end{enumerate}

  \begin{minted}[texcomments,linenos,escapeinside=||]{sql}
    /* L is a linked list */
    solve(L):
      if L is empty
        return 1
      else
        s = 1
        B = new linked list
        for j from 1 to 3 /* inclusive */
          i = 0
          A = L.head
          while A.next is not null
            b = (i >> j) & 3  /* $j$th and $(j+1)$th least significant bits of $i$ */
            if b == 1
              B.add(i / 2, A.value)
            A = A.next
            i = i + 1
          s = s + helper(B)
    \end{minted}

    \color{blue}
    \textbf{Answer:}
    \begin{enumerate}[(i)]
    \item $T(n) = 3 \cdot T(\frac{n}{4}) + \Theta(n )$
    \item $T(n) = \Theta(n)$
    \item Case 3.
    \item If there exists a constant $ \epsilon > 0 $ such that $ f(n) = \Omega(n^{\log_b a + \epsilon}) $, and if $ f(n) $ additionally satisfies the regularity condition $ a \cdot f\left(\frac{n}{b}\right) \leq c \cdot f(n) $ for some constant $ c < 1 $ and all sufficiently large $ n $, then $ T(n) = \Theta(f(n))$.
    
    \item $ 3 \cdot \frac{n}{4} \leq c \cdot n $, if $ c = \frac{3}{4} < 1 $ \xrightarrow[]{} $ n \leq n $
  \end{enumerate}
   \textbf{Justification:} \\
The loop iterates three times, matching the number of recursive calls.\\ The expression $b = (i >> j) \& 3$ adds every fourth element to list $B$ because the two bits extracted by $(i >> j) \& 3$ repeat cyclically every four numbers, and only the combination $01$ ($b == 1$) satisfies the condition. The function $f(n)$ is equal to $\Theta(n)$ because the reduction in the size of the list runs in linear time.\\
\\
For $T(n) = 3 \cdot T(\frac{n}{4}) + \Theta(n )$, where $ a = 3 > 0 $, $ b = 4 > 1 $, and $ f(n) = \Theta(n ) $, the recurrence relation satisfies the conditions of the Master Theorem:
    \begin{itemize}
        \item $ \log_{b} a =\log_{4} 3 \approx 0.793 $\\
        \item Case 1: \\
        \hspace{30mm} $ f(n) = \Theta(n ) = O(n^{\log_{b} a - \epsilon})= O(n^{0.793 - \epsilon}) $ - No, because $ \epsilon > 0 $ and $\Theta(n )$ grows faster.
        \item Case 2: \\
        \hspace{30mm} $ f(n) = \Theta(n )  = \Theta(n^{\log_{b} a}\cdot \log^k n)= \Theta(n^{0.793} \cdot \log^k n) $ - No, because $\Theta(n)$ grows faster than $\Theta(n^{0.793} \cdot \log^k n)$ for any $k$.
        \item Case 3: \\
        \hspace{30mm} $ f(n) = \Theta(n ) = \Omega(n^{\log_{b} a + \epsilon}) = \Omega(n^{0.793 + \epsilon}) $ - Yes. For $ \epsilon \approx 0.2 $, $ n^{0.793 + \epsilon} $ satisfies the condition. \\
        \hspace{30mm}[\cite{Cormen2022}, Theorem 4.1 (Master theorem)] Let's check the regularity condition $ a \cdot f\left(\frac{n}{b}\right) \leq c \cdot f(n) $: \\
        \hspace{35mm} $ 3 \cdot \Theta(\frac{n}{4} ) \leq c \cdot \Theta(n ) $ \\
        \hspace{35mm}$ 3 \cdot \frac{n}{4} \leq c \cdot n $, if $ c = \frac{3}{4} < 1 $ \\
        \hspace{35mm} $ n \leq n $. This condition is satisfied.
    \end{itemize}
    Answer: $ \Theta(n) $
    \color{black}

  \clearpage
  \item Consider the following modification of the binary search algorithm~\cite[Exercise~2.3-6]{Cormen2022}:
  given an input array $A$ of size $n$, a parameter $k$, and an integer value $x$
  \begin{enumerate}
    \item if the array is empty — stop with failure ($x$ not found)
    \item if the size of the input array is $1$ and $A[0] = x$ — stop with success ($x$ is found)
    \item otherwise, we create midpoints $m_i$ (for all integer indices $1 \leq i \leq k$)
      \emph{evenly} spaced between start index $m_0 = 0$ and end index $m_{k+1} = n - 1$
    \item we then find consecutive indices $m_i$ and $m_{i+1}$ such that $A[m_i] \leq x \leq A[m_{i+1}]$
      and repeat the search on the subarray $A[m_i:m_{i+1}]$
  \end{enumerate}

  Compute the time complexity of this algorithm in terms of $n$ and $k$ using two different methods:
  \begin{enumerate}[(i)]
    \item Write down the running time $T(n)$ as a recurrent formula.
    \item Apply the master theorem and get closed form formula for $T(n)$ (assume $k$ is a constant, the final answer may be independent of $k$):
      \begin{enumerate}
        \item Which case of the master theorem applies?
        \item Write down the conditions for this case that need to be checked (write down specialized version for this particular recurrence).
        \item Provide asymptotic complexity for $T(n)$ using $\Theta$-notation.
      \end{enumerate}

    \item Apply the substitution method and get closed form formula for $T(n, k)$ (treat $k$ as a parameter, the final answer should depend on $k$).
    \item Which value of $k$ is optimal? Justify your answer (4--6 sentences).
  \end{enumerate}

  \color{blue}
  \textbf{Answer:}
  \begin{enumerate}[(i)]
    \item $T(n) = T(\frac{n}{k+1}) + \Theta(k)$
    \item $T(n) = \Theta(\log n)$
    \item $T(n,k) = \Theta(k\cdot \log_{k+1} n)$
    \item Optimal value of $k$ equal to $1$
  \end{enumerate}
  \textbf{Justification: }
  \begin{enumerate}[(i)]
    \item The general form of the recurrence relation is as follows:
\[
T(n) = a \cdot T\left(\frac{n}{b}\right) + f(n)
\]
\hspace{20em} In the binary search algorithm:
\begin{itemize}
    \item We have $a = 1$, because the function is called only once to repeat the search on the subarray $A[m_i : m_{i+1}]$.
    \item $b = k + 1$, because the algorithm divides the array into $ k + 1 $ subarrays.
\end{itemize}

Thus, the recurrence relation for the algorithm is:
\[
T(n) = T\left(\frac{n}{k+1}\right) + \Theta(k)
\]



\item For $T(n) = T(\frac{n}{k+1}) + \Theta(k)$, where $ a = 1 > 0 $, $ b = k + 1 > 1 $ ($k \geq 1 $), and $ f(n) = \Theta(k)$, the recurrence relation satisfies the conditions of the Master Theorem:
    \begin{itemize}
        \item $ \log_{b} a =\log_{k+1} 1 = 0 $
        \item Case 1: \\
        \hspace{30mm} $ f(n) = \Theta(k) \neq O(n^{-\epsilon}) $ for any $ \epsilon > 0 $, because $ \Theta(k) $ is bounded below by a positive constant, while $ O(n^{-\epsilon}) $ tends to zero as $ n \to \infty $.
        \item Case 2: \\
        \hspace{30mm} $ f(n) = \Theta(k) = \Theta(n^{\log_{b} a}\cdot lg^k n) = \Theta(n^0\cdot lg^k n) = \Theta(lg^k n)$ - Yes, if $k = 0$.
    \end{itemize}
    Therefore, $\Theta(\log^1 n) = \Theta(\log n)$




    
   \item The substitution method:
    \begin{itemize}
        Suppose $n = b^i = (k+1)^i$\\
        \hspace{30mm} $T((k+1)^i) = T((k+1)^{i-1}) + \Theta(k)$\\
        \hspace{30mm} $T((k+1)^i) = T((k+1)^{i-2}) + \Theta(k) \cdot 2$\\
        \hspace{30mm} $T((k+1)^i) = T((k+1)^{i-3}) + \Theta(k) \cdot 3$\\
        \hspace{30mm} $\dots$\\
        \hspace{30mm} $T((k+1)^i) = T(1) + \Theta(k) \cdot i$\\
        \hspace{30mm} $i = \log_{k+1} n$, $T(n) = T(1) + \Theta(k) \cdot \log_{k+1} n$\\
        \hspace{30mm} $T(n) = \Theta(k \cdot \log_{k+1} n)$\\
    \end{itemize}

\item Optimal value of $k$ is $1$ because $k \geq 1$ and the base of the logarithm satisfies the condition since $k + 1 \geq 2 > 1$. Since $\Theta(k \cdot \log_{k+1} n)$ increases with larger values of $k$, it is advantageous to choose the smallest possible value.
  \end{enumerate}
  \color{black}

  \clearpage
  \item For each of the following recurrences, apply the master theorem~\cite[Theorem~4.1]{Cormen2022} yielding a closed form formula
  for the asymptotic complexity of $T(n)$.
  Assume that $T(n) = 1$ when $n < 10$ for all recurrences below. You must specify the following:
  \begin{itemize}
    \item Can the theorem be applied to the recurrence?
    \item If yes, then
    \begin{enumerate}[(i)]
      \item Which case of the master theorem applies?
      \item Which conditions for this case that need to be checked (write down specialized version for a particular recurrence).
      \item Prove that the condition is satisfied.
      For asymptotic notation, each property used in a proof \textbf{must} be either
      proven explicitly or properly referenced (e.g. by citing \cite{Cormen2022} or a particular Lecture/Tutorial slide).
      \item Provide asymptotic complexity for $T(n)$ using $\Theta$-notation.
    \end{enumerate}
    \item Otherwise, provide an explicit justification, explaining why the theorem cannot be applied.
  \end{itemize}
  \begin{enumerate}
    \item $T(n) = 6 T(n / 36) + \sqrt{n} \cdot \log_2^2 n$
    \item $T(n) = 4 T(3 n / 8) + n^2$
    \item $T(n) = 7 T(n / 5) + \log_2 \left(n!\right)$
    \item $T(n) = 3 T(\sqrt[3]{n}) + n$               
    \item $T(n) = \frac{1}{3} T(3n) + n \log_3 n$     
  \end{enumerate}

  \color{blue}
  \textbf{Answer:}
   \begin{enumerate}
    \item $T(n) = \Theta(\sqrt{n} \cdot lg^3n)$
    \item $T(n) =  \Theta(n^2)$
    \item $T(n) = \Theta (n^{log_57})$
    \item No answer              
    \item No answer    
  \end{enumerate}
  \textbf{Justification: }
    \begin{enumerate}

    
\item For $T(n) = 6 T(n / 36) + \sqrt{n} \cdot \log_2^2 n$, where $ a = 6 > 0 $, $ b = 36 > 1 $, and $ f(n) = \sqrt{n} \cdot \log_2^2 n $, the recurrence relation satisfies the conditions of the Master Theorem:
    \begin{itemize}
        \item $ \log_{b} a =\log_{36} 6 = \frac{1}{2} $
        \item Case 1: \\
        \hspace{30mm} $ f(n) = \sqrt{n} \cdot \log_2^2 n= O(n^{\log_{b} a - \epsilon}) = O(n^{\frac{1}{2} - \epsilon}) $ - No, because $ \epsilon > 0 $ and $ \sqrt{n} \cdot \log_2^2 n $ grows faster.
        \item Case 2: \\
        \hspace{30mm} $ f(n) = \sqrt{n} \cdot \log_2^2 n = \Theta(n^{\log_{b} a}\cdot lg^k n) = \Theta(n^{\frac{1}{2}}\cdot lg^k n) $ - Yes, if $k = 2$.
    \end{itemize}
    Answer: $\Theta(\sqrt{n} \cdot lg^{k+1}n) = \Theta(\sqrt{n} \cdot lg^3 n)$


    
\item For $ T(n) = 4 T\left(\frac{3n}{8}\right) + n^2 $, where $ a = 4 > 0 $, $ b = \frac{8}{3} > 1 $, and $ f(n) = n^2 $, the recurrence relation satisfies the conditions of the Master Theorem:
    \begin{itemize}
        \item $ \log_{b} a =\log_{\frac{8}{3}} 4 \approx 1.413 $\\
        \item Case 1: \\
        \hspace{30mm} $ f(n) = n^2 = O(n^{\log_{b} a - \epsilon})= O(n^{1.413 - \epsilon}) $ - No, because $ \epsilon > 0 $ and $ n^2 $ grows faster.
        \item Case 2: \\
        \hspace{30mm} $ f(n) = n^2  = \Theta(n^{\log_{b} a}\cdot \log^k n)= \Theta(n^{1.413} \cdot \log^k n) $ - No, because $ n^2 $ grows faster than $ n^{1.413} \cdot \log^k n $ for all $ k $.
        \item Case 3: \\
        \hspace{30mm} $ f(n) = n^2 = \Omega(n^{\log_{b} a + \epsilon}) = \Omega(n^{1.413 + \epsilon}) $ - Yes. For $ \epsilon \approx 0.6 $, $ n^{1.413 + \epsilon} $ satisfies the condition. \\
        \hspace{30mm} Let's check the regularity condition $ a \cdot f\left(\frac{n}{b}\right) \leq c \cdot f(n) $: \\
        \hspace{35mm} $ 4 \cdot \frac{9 \cdot n^2}{64} \leq c \cdot n^2 $ \\
        \hspace{35mm} $ \frac{9 \cdot n^2}{16} \leq c \cdot n^2 $, if $ c = \frac{9}{16} < 1 $ \\
        \hspace{35mm} $ n^2 \leq n^2 $. This condition is satisfied.
    \end{itemize}
    Answer: $ \Theta(n^2) $

\item For $T(n) = 7 T(n / 5) + \log_2 \left(n!\right)$, where $ a = 7 > 0 $, $ b = 5 > 1 $, and $ f(n) = \log_2 \left(n!\right) $, the recurrence relation satisfies the conditions of the Master Theorem:
    \begin{itemize}
        \item $\log_{b} a = \log_{5} 7 \approx 1.209 $
        \item Case 1: \\
        \hspace{30mm} $ f(n) = \log_2 \left(n!\right) = O(n^{\log_{b} a - \epsilon}) = O(n^{1.209 - \epsilon}) $ - Yes, because $\log_2 \left(n!\right)$ does not grow faster than $n^{1.209}$.
    \end{itemize}
    Answer: $\Theta(n^{\log_{5} 7})$



    
\item For $T(n) = 3 T(\sqrt[3]{n}) + n$, the recurrence relation does not satisfy the conditions of the Master Theorem, because $T(\sqrt[3]{n})$ does not match the form $T(n / b)$, where $b > 1$ is a constant. Therefore, the Master Theorem cannot be applied directly.


\item For $T(n) = \frac{1}{3} T(3n) + n \log_3 n$, the recurrence relation does not satisfy the conditions of the Master Theorem, because $b = \frac{1}{3}$ does not match $b>1$. Therefore, the Master Theorem cannot be applied directly.  
  \end{enumerate}
  \color{black}

  \clearpage
  \item \textbf{(+1\% extra credit)} Consider the following recurrence with parameters $a > 1$ and $b > 1$:
  \[
    T(n) = a T(n / b) + n^{4 \log_a b} \ln^{a - 2b}  n
  \]

  Specify \textbf{precise} conditions on $a$ and $b$ such that
  \begin{enumerate}[(i)]
    \item Case 1 of the master theorem applies to $T(n)$
    \item Case 2 of the master theorem applies to $T(n)$
    \item Case 3 of the master theorem applies to $T(n)$
    \item The master theorem is not applicable to $T(n)$
  \end{enumerate}

  Provide a brief justification in each case.

  \color{blue}
  \textbf{Answer:} \emph{(replace with your answer)} \\
  \textbf{Justification: } \emph{(replace with your justification)}
  \color{black}

\end{enumerate}

\begin{thebibliography}{BoasWrench}

\bibitem[CLRS]{Cormen2022}
Cormen, T.H., Leiserson, C.E., Rivest, R.L. and Stein, C., 2022. \emph{Introduction to algorithms, Fourth Edition}. MIT press.

\end{thebibliography}

\end{document}